\documentclass[12pt]{article}

\usepackage[T1]{fontenc}
\usepackage[utf8]{inputenc}
\usepackage[lining]{ebgaramond}                 % Garamond revival ~ the original old-style type
\usepackage{amsmath,amssymb}
\usepackage[italic,defaultmathsizes,noplusnominus]{mathastext} % math variables in the text italic
\usepackage{geometry}
\usepackage{fancyhdr}
\usepackage{pifont}                     % printer's ornaments
\usepackage{array}
\usepackage{graphicx}                   % \reflectbox for mirrored ornaments
\usepackage{lettrine}                   % opening drop capital
\renewcommand{\LettrineTextFont}{\normalfont}
\usepackage{microtype}
\usepackage{etoolbox}
\usepackage{tabularx}
\usepackage[perpage]{footmisc}
% The original sets its tables of residues a size down from the body text.
\AtBeginEnvironment{tabular}{\small}

%%% --- trim size of the original volume (543.84 x 706.8 pt) -------------
%  Measured off the scan of the original: text measure 360 pt, 32 lines to
%  the page, 15.60 pt baseline to baseline, text block 518 pt deep.
\geometry{
  paperwidth  = 543.84pt,
  paperheight = 706.80pt,
  left        = 91.92pt,
  right       = 91.92pt,
  top         = 72pt,
  textheight  = 558pt,
  headheight  = 14pt,
  headsep     = 9pt,
}

%%% --- printer's ornaments and rules -----------------------------------
\newcommand{\ornL}{\raisebox{-0.25ex}{\ding{167}}}
\newcommand{\ornR}{\raisebox{-0.25ex}{\reflectbox{\ding{167}}}}

% the running head of the original:  ornament ) 64 ( ornament
\newcommand{\pageornament}[1]{%
  {\normalsize\ornL\kern.45em\textrm{)}\kern.40em #1\kern.40em\textrm{(}\kern.45em\ornR}}

% heavy band under the running head on the opening page of E552
\newcommand{\thickband}{%
  \par\nointerlineskip\vspace{.4ex}%
  \noindent\rule{\textwidth}{2.4pt}\par\vspace{1.0ex}}

% floral band on the opening page of E554
\newcommand{\ornband}{%
  \par\vspace{.6ex}\noindent
  \makebox[\textwidth]{\ornL\hfill\ornR\hfill\ornL\hfill\ornR\hfill\ornL%
                       \hfill\ornR\hfill\ornL\hfill\ornR\hfill\ornL}%
  \par\vspace{1.0ex}}

% short centred rule closing a title block
\newcommand{\shortrule}{%
  \par\vspace{1.2ex}{\centering\rule{0.34\textwidth}{1.3pt}\par}\vspace{1.4ex}}

% wavy rule closing the title block of E554 (approximated by a double rule)
\newcommand{\wavyrule}{%
  \par\vspace{1.2ex}{\centering\rule{0.62\textwidth}{1.1pt}\\[-0.9ex]%
  \rule{0.62\textwidth}{0.5pt}\par}\vspace{1.6ex}}

% plain horizontal divider used between the numbered tables of E554
\newcommand{\tabrule}{\par\vspace{.5ex}\noindent\rule{\textwidth}{1.1pt}\par\vspace{.5ex}}

%%% --- title block of a memoir -----------------------------------------
\newcommand{\worktitle}[1]{{\centering\Large{#1}\par}\vspace{.45ex}}
\newcommand{\worttiny}[1]{{\centering\normalsize{#1}\par}\vspace{.45ex}}
\newcommand{\workbig}[1]{{\centering\LARGE{#1}\par}\vspace{.45ex}}
\newcommand{\worksub}[1]{{\centering\large{#1}\par}\vspace{.45ex}}

%%% --- centred section heading (Hypothesis. / Corollary 1. / Scholium.) -
\newcommand{\heading}[1]{%
  \par\addvspace{2.0ex plus .5ex minus .3ex}%
  {\centering\large #1\par}%
  \addvspace{1.0ex plus .3ex minus .2ex}}

%%% --- Euler's paragraph number:  §. 4. ---------------------------------
\newcommand{\pnum}[1]{\textbf{\S.\,#1.}\kern1.0em}

%%% --- original page boundary -------------------------------------------
%  Starts a new physical page and forces its number to the number printed
%  in the 1783 edition, so the running heads match the original throughout.
\newcommand{\origpage}[1]{\clearpage\setcounter{page}{#1}}

%%% --- running head / feet ----------------------------------------------
\pagestyle{fancy}
\fancyhf{}
\fancyhead[C]{\pageornament{\thepage}}
\renewcommand{\headrulewidth}{0pt}
\renewcommand{\footrulewidth}{0pt}
\fancypagestyle{plain}{\fancyhf{}\fancyhead[C]{\pageornament{\thepage}}%
  \renewcommand{\headrulewidth}{0pt}}

%%% --- paragraph shape --------------------------------------------------
\setlength{\parindent}{1.6em}
\setlength{\parskip}{0pt plus .3pt}
\linespread{1.076}   % 12pt type on the original's measured 15.60pt leading
\frenchspacing

%%% --- display maths spacing, close to the original's tight setting -----
\setlength{\abovedisplayskip}{1.1ex plus .3ex minus .2ex}
\setlength{\belowdisplayskip}{1.1ex plus .3ex minus .2ex}
\setlength{\abovedisplayshortskip}{0.7ex plus .2ex}
\setlength{\belowdisplayshortskip}{0.9ex plus .2ex}

%%% --- tables of residues ------------------------------------------------
\newenvironment{restable}[1]      % #1 = tabular column spec
  {\par\addvspace{1.0ex}\centering\begin{tabular}{#1}}
  {\end{tabular}\par\addvspace{1.0ex}}
\begin{document}

%%% ---- original page 121 (PDF page 1) ----
%%% signature line: "Euleri Opusc. Anal. Tom. I."  /  "Q"   catchword: "$\alpha \pm p$"
\origpage{1}

\vspace*{0cm}
\workbig{ON RAISING INDUCTION}
\worktitle{TO FULL CERTAINTY \footnote{Written in 1775 and originally published as \textit{De inductione ad plenam certitudinem evehenda}, Acta Academiae Scientiarum Imperialis Petropolitanae, Volume 1780: II, pp. 38-48. E566 in the Eneström index. Translated from the Latin by Samy Lahlou Kamal in October 2026, email: samy.lahloukamal@kellogg.ox.ac.uk.}}
\vspace{0.7cm}
\worksub{Author}
\worktitle{L. EULER}
\vspace{0.7cm}


\heading{\pnum{1}}

\lettrine[lines=2,findent=.15em,nindent=0pt]{I}{t} is well known that most properties of numbers were first observed through induction alone, and that geometers subsequently strove to confirm them with rigorous proofs; Fermat, in particular, devoted himself to this task with great zeal and considerable success. It is, however, deeply to be regretted that many of his proofs have been completely lost through the ravages of time, especially those concerning the resolution of numbers into three triangular numbers, four squares, five pentagonal numbers, six hexagonal numbers, and so forth, which he had asserted he had discovered. Subsequently, the resolution into four squares was indeed established by the renowned La Grange\footnote{\textit{Démonstration d'un théorème d'arithmétique}, Nouveaux mémoires de l'Académie royale des sciences et belles-lettres de Berlin, année 1770, p.189-201} and myself \footnote{\textit{Novae demonstrationes circa resolutionem numerorum in quadrata}, Nova acta eruditorum, Volume 1773, pp. 193-211, E445 in the Eneström index.}on the basis of absolutely certain principles; yet, this method of demonstration has not yet been successfully adapted to triangular and other polygonal numbers. Meanwhile, however, these resolutions have already been raised, solely through induction, to such a degree of certainty that there appears to be no room left for further doubt. Moreover, the process of induction itself can perhaps be so reinforced by certain arguments that it may be regarded as equivalent to an absolute demonstration, a point I shall attempt to illustrate in connection with the decomposition of numbers into four squares and three triangular numbers.\\

\heading{On the Resolution of Numbers into Four Squares.}

\pnum{2} Although this resolution is already fortified by a complete demonstra-tion, it will nevertheless be worthwhile to examine how far we can succeed in establishing its truth by relying solely on induction. Let us, therefore, first call to our aid the famous theorem of Fermat, that every prime number of the form $4n+1$ is always the sum of two squares, a proof of which I have, in fact, already long since presented\footnote{Translator: Euler's first proof starts in the article \textit{De numeris, qui sunt aggregata duorum quadratorum}, Novi Commentarii academiae scientiarum Petropolitanae, Volume 4, pp. 3-40, E228 in the Eneström index, and is completed in the article \textit{Demonstratio theorematis Fermatiani omnem numerum primum formae 4n+1 esse summam duorum quadratorum}, Novi Commentarii academiae scientiarum Petropolitanae, Volume 5, pp. 3-13, E241 in the Eneström index.}. From this, therefore, it immediately follows that the sum of two such numbers is certainly the sum of four squares. Moreover, the sums of two numbers of the form $4n+1$ yield numbers of the form $4n+2$, that is, numbers that are singly even. Consequently, if all numbers of the form $4n+2$ could be resolved into two prime numbers of the form $4n+1$, they would also be sums of four squares; and if this were demonstrated, then the halves of these numbers, that is, all odd numbers, and indeed absolutely all numbers, would likewise be sums of four squares. For it can easily be shown that if any number $N$ is the sum of four squares, then its double, its quadruple, and in general the form $2^{\pi}N$, will also always be the sum of four squares. Moreover, if $N$ is an even number, as in our case, namely $4n + 2$, then its half is also certainly the sum of four squares.\\

\pnum{3} The entire question, therefore, reduces to this: whether all numbers that are oddly even, that is, of the form $4n + 2$, can always be split into two parts of the form $4n + 1$ that are simultaneously prime numbers, a matter we shall investigate by means of induction. Having identified the prime numbers of the form $4n + 1$, namely 1, 5, 13, 17, 29, 37, 41, 53, 61, 73, etc., let us run through all numbers of the form $4n + 2$ in order, and for each one, append its resolution into two prime numbers of this kind:\\

\begin{tabular}{l|l|l|l} 
\parbox[t]{0.2\linewidth}{%
    \vspace*{-\baselineskip}
    \begin{align*} 
        2 &= \quad \! 1 + 1 \\ 
        6 &= \quad \! 1 + 5 \\
        10 &= \quad \! 5 + 5 \\
        14 &= \quad \! 1 + 13 \\
        18 &= \begin{cases}
           1 + 17 \\ 5 + 13
        \end{cases} \\
        22 &= \quad \! 5 + 17 \\
        26 &= \ \ \ \! 13 + 13 \\
        30 &= \begin{cases}
           1 \ + 29 \\ 13 + 17
        \end{cases} \\
        34 &= \begin{cases}
           5 + 29 \\ 17 + 17
        \end{cases} \\
        38 &= \quad \! 1 + 37 \\
        42 &= \begin{cases}
           1 + 41 \\ 5 + 37 \\ 13 + 29
        \end{cases} 
    \end{align*}} &  
\parbox[t]{0.2\linewidth}{%
    \vspace*{-\baselineskip}
    \begin{align*} 
        46 &= \begin{cases}
           5 + 41 \\ 17 + 29
        \end{cases} \\
        50 &= \quad \! 13 + 37 \\
        54 &= \begin{cases}
           1 + 53 \\ 13 + 41 \\ 17 + 37
        \end{cases} \\
        58 &= \begin{cases}
           5 + 53 \\ 17 + 41 \\ 29 + 29
        \end{cases} \\
        62 &= \quad \! 1 + 61 \\
        66 &= \begin{cases}
           5 + 61 \\ 13 + 53 \\ 29 + 27
        \end{cases} \\
        70 &= \begin{cases}
           17 + 53 \\ 29 + 41
        \end{cases} 
    \end{align*}} &
\parbox[t]{0.2\linewidth}{%
    \vspace*{-\baselineskip}
    \begin{align*} 
        74 &= \begin{cases}
           1 + 73 \\ 13 + 61 \\ 37 + 37
        \end{cases} \\
        78 &= \begin{cases}
           5 + 73 \\ 17 + 61 \\ 37 + 41
        \end{cases} \\
        82 &= \begin{cases}
           29 + 53 \\ 41 + 41
        \end{cases} \\
        86 &= \quad \! 13 + 76 \\
        90 &= \begin{cases}
           1 + 89 \\ 17 + 73 \\ 29 + 61 \\ 37 + 53
        \end{cases} \\
        94 &= \begin{cases}
           5 + 89 \\ 41 + 53
        \end{cases} 
    \end{align*}} &  
\parbox[t]{0.2\linewidth}{%
    \vspace*{-\baselineskip}
    \begin{align*} 
        98 &= \begin{cases}
           1 + 97 \\ 37 + 61
        \end{cases} \\
        102 &= \begin{cases}
           1 + 101 \\ 5 + 97 \\ 13 + 89 \\ 29 + 73 \\ 41 + 61
        \end{cases} \\
        106 &= \begin{cases}
           5 + 101 \\ 17 + 89 \\ 53 + 53
        \end{cases} \\
        110 &= \begin{cases}
           1 + 109 \\ 13 + 97 \\ 37 + 73
        \end{cases}
    \end{align*}}
\end{tabular}

\vspace{1cm}

\pnum{4} Such a resolution, therefore, succeeded successfully up to the number 110; and if anyone wishes to continue this task, he will discover that it not only always succeeds but also that further resolutions continually present themselves. At the same time, however, it would be a rash conclusion to infer from this that all numbers of the form $4n + 2$ are invariably resolvable into two [prime] numbers of the form $4n + 1$. Although multiple resolutions generally exist, there are nevertheless cases where only a single resolution is found, such as with the numbers 26, 38, 50, 62, and 86. From this, one might object that not only do such cases occur among larger numbers, but it might even happen that such a resolution fails entirely, though this seems less likely to occur the further we proceed. However, we must also acknowledge that prime numbers of the form $4n + 1$ become increasingly rare, such that, for this reason, many resolutions are excluded. Moreover, if prime numbers of the form $4n + 1$ were to cease entirely among the largest numbers, then such a resolution could no longer be possible. But since Euclid has already demonstrated that the set of prime numbers is indeed infinite, this must also hold true for prime numbers of the form $4n + 1$, which therefore can never run out completely.\\

\pnum{5} Therefore, for these reasons, it is not appropriate to insist on the assertion that all numbers of the form $4n + 2$ admit a resolution into two primes of the form 
$4n + 1$, even though this may perhaps be true and could conceivably be confirmed at some point by a rigorous proof. Meanwhile, however, it is evident that if the number of primes of the form $4n + 1$ employed in these resolutions were greater than what we have assumed here, there would be far less reason to fear cases that might resist such a resolution. To those numbers, moreover, we may quite rightly add all square numbers of the form 4n + 1; for when added to prime numbers of the form 4n + 1, they yield numbers that are resolvable into three squares. Therefore, let us also add odd squares to those numbers of the form 4n + 1; in this way, up to the number 100, we obtain the following numbers for our resolutions:
$$1, \ \ 5, \ \ 9, \ \ 13, \ \ 17, \ \ 25, \ \ 29, \ \ 37, \ \ 41, \ \ 49, \ \ 53, \ \ 61, \ \ 73, \ \ 81, \ \ 89, \ \ [97].$$

\pnum{6} Let us therefore see how the resolutions of numbers of the form $4n+2$
into two such parts proceed. \\


\begin{tabular}{l|l|l|l} 
\parbox[t]{0.2\linewidth}{%
  \vspace*{-\baselineskip}
    \begin{align*} 
        2 &= \quad \! 1 + 1 \\ 
        6 &= \quad \! 1 + 5 \\
        10 &= \begin{cases}
          1 + 9 \\ 5 + 5
        \end{cases} \\
        14 &= \begin{cases}
          1 + 13 \\ 5 + 9
        \end{cases} \\
        18 &= \begin{cases}
           1 + 17 \\ 5 + 13 \\ 9 + 9
        \end{cases} \\
        22 &= \begin{cases}
           5 + 17 \\ 9 + 13
        \end{cases} \\
        26 &= \begin{cases}
          1 + 25 \\ 9 + 17 \\ 13 + 13
        \end{cases} \\
        30 &= \begin{cases}
           1 + 29 \\ 5 + 25 \\ 13 + 17
        \end{cases} \\
        34 &= \begin{cases}
           5 + 29 \\ 9 + 25 \\ 17 + 17
        \end{cases} \\
        38 &= \begin{cases}
           1 + 37 \\ 9 + 29 \\ 13 + 25
        \end{cases} 
    \end{align*}} &  
\parbox[t]{0.2\linewidth}{%
    \vspace*{-\baselineskip}
    \begin{align*} 
      42 &= \begin{cases}
        1 + 41 \\5 + 37 \\ 13 + 29 \\ 17 + 25
      \end{cases}\\
      46 &= \begin{cases}
           5 + 41 \\ 9 + 37 \\ 17 + 29
        \end{cases} \\
      50 &= \begin{cases}
        1 + 49 \\ 9 + 41 \\ 13 + 37 \\ 25 + 25
      \end{cases} \\
      54 &= \begin{cases}
           1 + 53 \\ 5 + 49 \\ 13 + 41 \\ 17 + 37 \\ 25 + 29
        \end{cases} \\
      58 &= \begin{cases}
           5 + 53 \\ 9 + 49 \\ 17 + 41 \\ 29 + 29
        \end{cases}
    \end{align*}} &
\parbox[t]{0.2\linewidth}{%
    \vspace*{-\baselineskip}
    \begin{align*} 
      62 &= \begin{cases}
        1 + 61 \\ 9 + 53 \\ 13 + 49 \\ 25 + 37
      \end{cases} \\
      66 &= \begin{cases}
        5 + 61 \\ 13 + 53 \\ 17 + 49 \\ 25 + 41 \\ 29 + 37
      \end{cases} \\
      70 &= \begin{cases}
        9 + 61 \\ 17 + 53 \\ 29 + 41
      \end{cases} \\
      74 &= \begin{cases}
        1 + 73 \\ 13 + 61 \\ 25 + 49 \\ 37 + 37
      \end{cases} \\
      78 &= \begin{cases}
        5 + 73 \\ 17 + 61 \\ 25 + 53 \\ 29 + 49 \\ 37 + 41
      \end{cases}
    \end{align*}} &  
\parbox[t]{0.2\linewidth}{%
    \vspace*{-\baselineskip}
    \begin{align*} 
      82 &= \begin{cases}
        1 + 81 \\ 29 + 53 \\ 41 + 41 \\ 9 + 73
      \end{cases} \\
      86 &= \begin{cases}
        5 + 81 \\ 13 + 73 \\ 26 + 61 \\ 37 + 49
      \end{cases} \\
      90 &= \begin{cases}
        1 + 89 \\ 9 + 81 \\ 17 + 73 \\ 29 + 61 \\ 37 + 53 \\ 41 + 49
      \end{cases} \\ 
      94 &= \begin{cases}
        5 + 89 \\ 13 + 81 \\ 41 + 53
      \end{cases} \\
      98 &= \begin{cases}
        1 + 97 \\ 9 + 89 \\ 17 + 81 \\ 25 + 73 \\ 37 + 61 \\ 49 + 49
      \end{cases}
    \end{align*}}
\end{tabular}

\vspace{0.3cm}

\pnum{7} Therefore, as squares have been added to our prime numbers of the form $4n + 1$, the number of possible resolutions has increased so remarkably, and will evidently increase much more with larger numbers, that there seems to be no reason to fear that there will ever be only a single resolution, let alone none at all; in this way, the induction performed earlier has been so strongly reinforced that it appears tantamount to a rigorous proof. Yet, should anyone still wish to object that the matter has not been sufficiently established, he would nevertheless be compelled to admit that, were it possible to further increase the number of such primes of the form 4n + 1, there would clearly be no room left for any further objection. \\

\pnum{8} Since those numbers of the form $4n + 1$ are suited to our purpose insofar as they are either squares themselves or resolvable into two squares, we may also include with them products formed from two or more numbers of the same type; for it has been demonstrated that the product of two or more numbers, each of which is the sum of two squares, can always be resolved into two squares. Therefore, to the aforementioned class of these numbers, we may further add the following: 
$$45, \ \ 65,\ \ 85,\ \ 117,\ \ 125,\ \ 145,\ \ 153,\ \ 185,\ \ 205,\ \ 221\ \ [etc.].$$
However, with the addition of these numbers, it is evident that the total number of resolutions increases so greatly, and continues to grow, that it is utterly impossible for such a resolution ever to fail. Thus, whereas in the preceding resolutions the number 82 yielded only four resolutions, the following ones will now apply:
$$82=1+ 81, \ \ 9 + 73,\ \ 17 +65,\ \ 29 + 53,\ \ 37 + 45,\ \ 41 + 41,$$
and thus, two new solutions were added. Therefore, on the basis of these arguments, the truth of the theorem, that every number is the sum of four squares, must be considered placed entirely beyond doubt. At the same time, however, we readily concede that this proof does not possess the degree of self-evidence usually required in rigorous demonstrations. \\

\heading{On the Resolution of Numbers into Three Triangular Numbers.}

\pnum{9\footnote{Translator: In the original article, this paragraph is numbered §5 instead of §9, the next one is numbered §6 instead of §10, etc. I chose to correct this mistake (numbering mistakes for paragraphs happened often in this period).}} If the number $N$ is the sum of three triangular numbers, that is
$$N = \frac{xx + x}{2} + \frac{yy + y}{2} + \frac{zz + z}{2},$$
it will be
$$8N+3 = (2x+1)^2 + (2y +1 )^2 + (2z+1)^2.$$
If, therefore, it can be demonstrated that all numbers of the form $8N + 3$ are sums of three squares, then it will also be demonstrated that all numbers N are sums of three triangular numbers; for a number of the form $8N + 3$ can only be equal to the sum of three odd squares. Therefore, if
$$8N+3 = (2x+1)^2 + (2y +1 )^2 + (2z+1)^2,$$
then conversely
$$N = \frac{xx + x}{2} + \frac{yy + y}{2} + \frac{zz + z}{2},$$
and consequently it is the sum of three triangular numbers. \\

\pnum{10} Now let us call to our aid the theorem mentioned above, stating that every prime number of the form $4n + 1$ is the sum of two squares, and thus $4n + 1 = aa + bb$; it is also evident that its double, $8n + 2$, will be the sum of two squares, for in this case
$$8n+2 = (a+b)^2 + (a-b)^2.$$
If, therefore, we add odd squares, which are of the form $8m + 1$, to each of these numbers of the form $8n + 2$, the resulting sums will certainly be sums of three squares; the form of these sums will be $8n + 3$. Consequently, the question reduces to this: can all numbers of the form $8N + 3$ be divided into two parts such that one is an odd square and the other is the double of a prime number of the form $4n + 1$? To investigate this at least for smaller numbers, let us list these doubles in order up to 200:
$$2, \ \ 10,\ \ 26,\ \ 34,\ \ 58,\ \ 74,\ \ 82,\ \ 106,\ \ 122,\ \ 146,\ \ 178,\ \ 194,\ \ 202.$$

\pnum{11} Let us now, therefore, run through all numbers of the form $8N+3$ and see whether they can all be composed from some number of this same class and an odd square, specifically, the sequence:
$$1,\ \ \ 9, \ \ \ 25,\ \ \ 49,\ \ \ 81,\ \ \ 121,\ \ \ 169, \ \ \ 225,$$
an investigation that may be presented as follows:\\

\begin{tabular}{l|l|l|l} 
\parbox[t]{0.2\linewidth}{%
  \vspace*{-\baselineskip}
    \begin{align*} 
        3 &= \quad \! 1 + 2 \\ 
        11 &= \begin{cases}
          1 + 10 \\ 9 + 2
        \end{cases} \\
        19 &= \quad \! 9 + 10 \\
        27 &= \begin{cases}
          1 + 26 \\ 25 + 2
        \end{cases} \\
        35 &= \begin{cases}
          1 + 34 \\ 9 + 26 \\ 25 + 10
        \end{cases} \\
        43 &= \quad \! 9+34 \\
        51 &= \begin{cases}
          25 + 26 \\ 49 + 2
        \end{cases} \\
        59 &= \begin{cases}
          1 + 58 \\ 25 + 34 \\ 49 + 10
        \end{cases} \\
        67 &= \quad \! 9 + 58
    \end{align*}} &  
\parbox[t]{0.2\linewidth}{%
  \vspace*{-\baselineskip}
    \begin{align*} 
        75 &= \begin{cases}
          1 + 74 \\ 49 + 26
        \end{cases} \\
        83 &= \begin{cases}
          1 + 82 \\ 9 + 74 \\ 25 + 58 \\ 49 + 34 \\ 81 + 2
        \end{cases} \\
        91 &= \begin{cases}
          9 + 82 \\ 81 + 10
        \end{cases} \\
        99 &= \quad \! 25 + 74 \\
        107 &= \begin{cases}
          1 + 106 \\ 25 + 82 \\ 49 + 58 \\ 81 + 26
        \end{cases} \\
        115 &= \begin{cases}
          9 + 106 \\ 81 + 34
        \end{cases}
    \end{align*}} &
\parbox[t]{0.2\linewidth}{%
  \vspace*{-\baselineskip}
    \begin{align*} 
        123 &= \begin{cases}
          1 + 122 \\ 49 + 74 \\ 121 + 2 
        \end{cases} \\
        131 &= \begin{cases}
          9 + 122 \\ 25 + 106 \\ 49 + 82 \\ 121 + 10
        \end{cases} \\
        139 &= \quad \! 81 + 58 \\
        147 &= \begin{cases}
          1 + 146 \\ 25 + 122 \\ 121 + 26
        \end{cases} \\
        155 &= \begin{cases}
          9 + 146 \\ 49 + 106 \\ 81 + 74 \\ 121 + 34
        \end{cases}
    \end{align*}} &  
\parbox[t]{0.2\linewidth}{%
  \vspace*{-\baselineskip}
    \begin{align*} 
        163 &= \quad \! 81 + 82 \\
        171 &= \begin{cases}
          25 + 146 \\ 49 + 122 \\ 169 + 2
        \end{cases} \\
        179 &= \begin{cases}
          1 + 178 \\ 121 + 58 \\ 169 + 10
        \end{cases} \\
        187 &= \begin{cases}
          9 + 178 \\ 81 + 106 
        \end{cases} \\
        195 &= \begin{cases}
          1 + 194 \\ 169 + 26 \\ 49 + 146 \\ 121 + 74
        \end{cases}
    \end{align*}}
\end{tabular}

\vspace{0.3cm}

\pnum{12} From this, therefore, it seems one might conclude without excessive rashness that all numbers of the form $8n + 3$ can be resolved into two parts, one of which is a square and the other the double of a prime number of the form $4n + 1$. For since the number of available squares increases continuously, the number of combinations also grows ever larger; hence, if such a resolution were to fail anywhere, it ought to have failed among the smaller numbers. Yet, since cases have arisen that admit only a single resolution, as happened with the numbers 67, 99, 139, and 163, and since such cases will undoubtedly occur again, one may reasonably fear that such a resolution might fail completely in some instance. Therefore, although this conclusion may well be true, let us not place our confidence in it. \\

\pnum{13} However, since that class of numbers of the form $8n + 2$ admits of considerable expansion, specifically by adding to them doubled odd squares which are themselves sums of two squares, there will be absolutely no further risk of the resolution ever becoming impossible. Furthermore, we can also add the doubles of those composite numbers of the form $4n + 1$, namely, those that are products of two or more prime numbers of the same form, whereby the following progression of numbers of the form $8n + 2$ will result:
\begin{align*}
  2, \ \ 10, \ \ &18, \ \ 26, \ \ 34, \ \ 50, \ \ 58, \ \ 74, \ \ 82, \ \ 90, \ \ 98, \ \ 106, \ \ 122, \\
  &130, \ \ 146, \ \ 162, \ \ 170, \ \ 178, \ \ 194, \ \ 202, \ \ \text{etc.} \\
\end{align*}

\pnum{14} Therefore, by calling upon these new numbers as an aid, let us resolve only those numbers that admitted a unique resolution in the previous table, in this way:
\begin{align*}
  19 &= 1 \ \ \! \! \ \! + \ \ \, 18 = \ \ \ \ \! \! 9 + 10 \\
  43 &= 9 + \ \  34= 25 + 18 \\
  67 &= 9 + \ \ 58 = \ \! 49 + 18 \\
  99 &= 1 \ \ \! \! \ \! + \ \ 98 = \ \ \ \ \! \! \ \! \ \! 9 + 90 = 25 + 74 = \ \ \! \ \! 49 + 50 = 81 + 18 \\
  139 &= 9 + 130 =  \ \! 49 + 90 = 81 \ \! \ \! \ \!  + 58 = 121 + 18 \\
  163 &= 1 \ \ \! \! \ \! + \ \! 162 = 81 \ \! \! \ \! \ \ \! + 82
\end{align*}
Here, then, we see that at least one solution has been obtained; yet with larger numbers, still more will have to be found. Having duly weighed these matters, let geometers judge whether this induction, corroborated by so many arguments, should not be considered equivalent to a solid demonstration. \\

\pnum{15} Although this conclusion appears to require no further confirmation, substantial support can nevertheless be drawn from a completely different principle, namely, that all prime numbers of the form $8n + 3$ are always of the form $2aa+bb$ and are therefore themselves resolvable into three squares, a fact that has already been rigorously demonstrated\footnote{\textit{Demonstrationes circa residua ex divisione potestatum per numeros primos resultantia}, Novi Commentarii academiae scientiarum Petropolitanae, Volume 18, pp. 85-135, E449 in the Eneström index.}. From this, it also follows that these same prime numbers, when increased by eight, are likewise sums of three squares. For if $8n + 3 = 2aa+bb$, then
$$8n+11 = (a+2)^2 + (a-2)^2 + bb,$$
consequently, we are certain that the aforementioned resolution can never fail for prime numbers of the form $8n + 3$, nor even for the subsequent ones of the form $8n + 11$; and this holds true for numbers of the form $8n + 35$ as well, since
$$8n+35 = (a+4)^2 + (a-4)^2 + bb,$$
and will therefore also apply to all numbers of the form $8n + 3 + 8pp$, where $p$ denotes any number whatsoever. Furthermore, since it has also been demonstrated that all prime numbers of the form $8n + 1$ are likewise numbers of the form $2aa + bb$, that such products of two prime numbers $(8p + 1)(8q + 3)$, which are
certainly numbers of the form $8n + 3$, are always numbers of the same form $2aa + bb$,
it follows that these numbers, whether increased by 8, by 32, or in general by
$8rr$, are certainly sums of three squares. With these considerations added,
Fermat's theorem on triangular numbers must be regarded as having been
demonstrated with the utmost rigor. Nevertheless, we are compelled to admit
that, in this way, the necessity of why this must be so cannot be clearly perceived.\\

\par\vspace{3.0ex}
{\centering\rule{0.56\textwidth}{2.6pt}\par}
\par\vspace{1.5ex}

\end{document}
